% Ad Atticum 6.9 (ad-atticum-06-09) --- English translation
% Translated by Claude (Anthropic) from the Latin (Perseus).
% Style guide: STYLE.md.

\ciceroLetterOpener{CICERO ATTICO salutem}

\ciceroSection{1} When I had landed at the Piraeus on the day before the Ides of October, I received your letter at once from my slave Acastus. Long though I had been waiting for it, I was surprised --- when I saw the sealed packet --- at its shortness, and when I opened it, at the further \textit{[Greek: confusion]} of the little letters; for yours are accustomed to be the most carefully composed and the clearest of any. To be brief: I gathered from what you had written that you had reached Rome on the twelfth day before the Kalends of October with a fever. Struck violently --- and no more than was my duty --- I at once asked Acastus. He told me that both to you and to him things seemed right enough, and that he had heard the same at home from your people: that there was nothing the matter to speak of. The fact that, at the close of the letter, you had written you had a slight touch of fever at the time of writing, seemed to bear that out. But I loved you all the same, and was amazed that, even so, you had written to me with your own hand. Of this, then, enough. For I trust --- such are your good sense and self-discipline --- and, by Hercules, as Acastus bids me, I am confident that by now you are as well as we wish.

\ciceroSection{2} That you received my letter through Turranius I am glad. \textit{[Greek: Be on your guard]}, if you love me, against \textit{[Greek: the self-promotion of the schemer]}: \textit{[Greek: exactly so]}. As for this --- which by Hercules is a great grief to me, for I loved the man --- see to it, however small the sum, that of the Precius legacy he shall touch absolutely nothing. You will say that I need the money for the apparatus of a triumph. In which matter, as you direct, you will find me neither \textit{[Greek: empty]} in the seeking nor \textit{[Greek: unduly puffed up]} in the renouncing.

\ciceroSection{3} I gathered from your letter that you had heard from Turranius that I had handed over the province to my brother. Did I not catch on, then, to the prudence of your letter? You wrote that you were \textit{[Greek: suspending judgement]}. What was there worth doubting, if there was anything at all to make leaving my brother behind --- and such a brother --- acceptable? That \textit{[Greek: suspension]} of yours seemed to me \textit{[Greek: a rejection]}, not \textit{[Greek: a withholding]}. You warned me about young Quintus Cicero, that on no account was I to leave him behind. \textit{[Greek: That was my own dream]}. We saw all the same things, as though we had talked them over together. It could not be done otherwise; and your \textit{[Greek: long-drawn]} doubt freed me from doubt of my own. But I think you have already received a more careful letter from me on this matter.

\ciceroSection{4} I was on the point of sending letter-carriers off to you tomorrow; whom I think will get there before our friend Saufeius. But for him to come to you without a letter from me was hardly proper. You, as you promise, will write me a full account of my little Tullia --- I mean, of Dolabella --- and of the commonwealth, which I foresee in the highest dangers, and of the censors, and especially of the statues and pictures, as to what is to be done, whether the question is to be brought forward. I gave this letter on the Ides of October, on which day, as you write, Caesar reaches Placentia with four legions. I ask you, what is to become of us? My place of station now is the citadel at Athens.
